\begin{abstract}
Subtitles carry speech, but a deaf or hard-of-hearing viewer still misses the other sounds of a video: a siren
behind the camera, a dog barking in the next room. Captioning practice leaves out any sound whose source can be
seen, so the sounds that matter most are the ones whose source is \emph{off screen}. This project builds a system
that finds those sounds and shows a generated picture of each one beside the video while it is heard.

Nothing is trained. The system chains ready-made open models: two sound detectors find candidate sounds, two
audio-language models (``listeners'', asked for example ``is a siren present?'') confirm weak ones, a
vision--language model checks whether the source is visible on screen, and a text-to-image model draws the
picture. Every rule that joins them was fixed on a development set before it was tested.

Scored per sound against human labels, the final system finds 24 of 65 needed off-screen sounds on a test set of
88 clips with 24 wrong pictures, a viewer cost of 2.41 per clip (a miss costs 4, a wrong picture 2). This is
significantly cheaper than showing nothing (2.96; difference $-0.55$, 95\,\% interval $[-1.09, -0.02]$, $p = 0.047$)
and cheaper than the September version of the system (2.91), which on the development set was worse than showing
nothing. The on-screen check is the reason: on the September test set it removed about half a wrong picture per clip
compared with drawing every sound ($p < 0.001$, Holm-corrected).

Two limits remain. About one needed sound in five cannot be found even with perfect decisions, because the detectors
barely hear it under speech or music; and the vision
model confuses ``the object is in the frame'' with ``the object is making this sound now''. The test set was read after
each accepted change, so the test numbers are a report on a seen set, and no study with deaf viewers was run.
\end{abstract}
